\subsection{ক্রায়োসার্জারি}

\begin{learningbox}
এই অংশ শেষে শিক্ষার্থী পারবে:
\begin{itemize}
\item ক্রায়োসার্জারির ধারণা ব্যাখ্যা করতে।
\item ক্রায়োসার্জারিতে ব্যবহৃত নিম্ন তাপমাত্রা ও cryogen সম্পর্কে বলতে।
\item চিকিৎসায় ক্রায়োসার্জারির ব্যবহার, সুবিধা ও সীমাবদ্ধতা লিখতে।
\item ICT কীভাবে আধুনিক চিকিৎসা প্রযুক্তিকে সহায়তা করে তা ব্যাখ্যা করতে।
\end{itemize}
\end{learningbox}

\subsubsection{ক্রায়োসার্জারির ধারণা}

Cryo শব্দের অর্থ শীতল বা অত্যন্ত ঠান্ডা। ক্রায়োসার্জারি হলো এমন চিকিৎসা পদ্ধতি যেখানে অত্যন্ত নিম্ন তাপমাত্রা ব্যবহার করে অস্বাভাবিক বা রোগাক্রান্ত tissue ধ্বংস করা হয়। এটি সাধারণত wart, skin lesion, cervical abnormal cell এবং কিছু tumor চিকিৎসায় ব্যবহৃত হয়।

\begin{definitionbox}{সংজ্ঞা}
ক্রায়োসার্জারি হলো চিকিৎসাবিজ্ঞানের এমন পদ্ধতি, যেখানে liquid nitrogen, argon gas বা অন্য cryogen ব্যবহার করে অস্বাভাবিক tissue জমিয়ে ধ্বংস করা হয়।
\end{definitionbox}

National Cancer Institute-এর cancer treatment dictionary অনুযায়ী cryosurgery-তে extreme cold ব্যবহার করে abnormal tissue ধ্বংস করা হয়। আধুনিক যন্ত্রে temperature control, imaging, probe placement ও monitoring-এ ICT গুরুত্বপূর্ণ ভূমিকা রাখে।

\subsubsection{ক্রায়োসার্জারিতে ব্যবহৃত উপাদান}

\begin{keypointbox}{ব্যবহৃত উপাদান}
\begin{itemize}
\item \textbf{Liquid nitrogen}: খুব নিম্ন তাপমাত্রা তৈরি করতে ব্যবহৃত।
\item \textbf{Argon gas}: cryoprobe-এ controlled freezing তৈরি করতে ব্যবহৃত হতে পারে।
\item \textbf{Cryoprobe}: নির্দিষ্ট tissue-তে ঠান্ডা প্রয়োগের যন্ত্র।
\item \textbf{Imaging system}: ultrasound, CT বা MRI দিয়ে target area দেখা।
\item \textbf{Monitoring system}: temperature ও treatment progress পর্যবেক্ষণ।
\end{itemize}
\end{keypointbox}

\subsubsection{ক্রায়োসার্জারির পদ্ধতি}

চিকিৎসক প্রথমে রোগীর problem area নির্ধারণ করেন। এরপর cryoprobe বা spray-এর মাধ্যমে target tissue-তে extreme cold প্রয়োগ করা হয়। ঠান্ডায় cell-এর ভেতরের পানি বরফে পরিণত হয়, cell membrane ক্ষতিগ্রস্ত হয় এবং অস্বাভাবিক tissue ধীরে ধীরে ধ্বংস হয়। কিছু ক্ষেত্রে freeze-thaw cycle একাধিকবার করা হয়।

\begin{examplebox}{সহজ উদাহরণ}
ত্বকের wart চিকিৎসায় ডাক্তার liquid nitrogen ব্যবহার করে wart-এর tissue জমিয়ে দেন। কয়েকদিন পরে ক্ষতিগ্রস্ত tissue শুকিয়ে পড়ে যেতে পারে।
\end{examplebox}

\subsubsection{ক্রায়োসার্জারির ব্যবহারক্ষেত্র}

\begin{itemize}
\item ত্বকের wart, mole বা কিছু lesion।
\item cervical pre-cancerous cell।
\item prostate, liver বা kidney tumor-এর কিছু ক্ষেত্রে image-guided cryoablation।
\item চোখ, dermatology এবং gynecology-র কিছু চিকিৎসা।
\end{itemize}

\begin{boardbox}{বোর্ড ফোকাস}
ক্রায়োসার্জারি লিখতে “ঠান্ডা দিয়ে operation” বললেই হবে না। উত্তর হবে: extreme cold, cryogen, abnormal tissue destruction, cryoprobe, চিকিৎসা ব্যবহার।
\end{boardbox}

\subsubsection{ক্রায়োসার্জারির সুবিধা}

\begin{keypointbox}{সুবিধা}
\begin{itemize}
\item অনেক ক্ষেত্রে incision বা বড় operation লাগে না।
\item bleeding ও infection risk তুলনামূলক কম হতে পারে।
\item recovery time কম হতে পারে।
\item নির্দিষ্ট target tissue-তে treatment দেওয়া যায়।
\item imaging ও computer monitoring ব্যবহার করলে precision বাড়ে।
\end{itemize}
\end{keypointbox}

\subsubsection{ক্রায়োসার্জারির অসুবিধা}

\begin{itemize}
\item সব রোগ বা tumor-এর জন্য উপযুক্ত নয়।
\item ভুল target হলে healthy tissue ক্ষতিগ্রস্ত হতে পারে।
\item pain, swelling, blister বা scar হতে পারে।
\item internal treatment-এ advanced imaging ও skilled doctor দরকার।
\item কখনো repeat treatment দরকার হতে পারে।
\end{itemize}

\begin{mistakebox}{সতর্কতা}
ক্রায়োসার্জারি নিজে নিজে করার চিকিৎসা নয়। Liquid nitrogen বা freezing chemical ভুলভাবে ব্যবহার করলে ত্বক ও tissue মারাত্মকভাবে ক্ষতিগ্রস্ত হতে পারে।
\end{mistakebox}

\begin{practicebox}
\textbf{MCQ}
\begin{enumerate}
\item ক্রায়োসার্জারিতে সাধারণত কী ব্যবহার করা হয়?\\
ক. অত্যন্ত উচ্চ তাপ \quad খ. অত্যন্ত নিম্ন তাপ \quad গ. শুধু laser light \quad ঘ. সাধারণ পানি
\item Cryoprobe-এর কাজ কী?\\
ক. virus scan \quad খ. target tissue-তে ঠান্ডা প্রয়োগ \quad গ. web browsing \quad ঘ. data backup
\end{enumerate}

\textbf{সংক্ষিপ্ত প্রশ্ন}
\begin{enumerate}
\item ক্রায়োসার্জারি কী?
\item ক্রায়োসার্জারির দুটি সুবিধা লেখ।
\end{enumerate}
\end{practicebox}

